\Article{13}
Als letztes Beispiel wollen wir die allgemeine Abbildung der Oberfläche eines Rotationsellipsoids auf der Kugeloberfläche betrachten. Für das Ellipsoid behalten wir die Bezeichnungen aus dem vorigen Artikel bei. Der Radius der Kugel sei \(=A\), und wir setzen
\[
\begin{aligned}
X&=A\cos T\sin U\\
Y&=A\sin T\sin U\\
Z&=A\cos U
\end{aligned}
\]
\SourceLeaf{213}%
\noindent Wendet man hier die allgemeine Lösung aus Artikel 5 an, so findet man für eine beliebige Funktion \(f\), dass man \(T\) dem Realteil und \(i\log\mathrm{cotang}\tfrac{1}{2}U\) dem Imaginärteil von
\[
f(t+i\log\{\mathrm{cotang}\tfrac{1}{2}w\mathord{.}(\frac{1-\varepsilon\cos w}{1+\varepsilon\cos w})^{\frac{1}{2}\varepsilon}\})
\]
\noindent gleichsetzen muss\footnote{Wir übergehen hier sowohl die zweite Lösung aus Artikel 5, die sich von der obigen nur durch den Austausch von \(-T\) gegen \(+T\) unterscheiden und einer seitenverkehrten Abbildung entsprechen würde, als auch den Fall eines länglichen Ellipsoids. Dessen Behandlung ergibt sich nach dem, was im vorigen Artikel dargelegt wurde, unmittelbar aus der Behandlung des abgeplatteten Ellipsoids.}.

Die einfachste Lösung erhält man, indem man \(f\upsilon=\upsilon\) setzt. Daraus folgt
\[
T=t,\qquad \mathrm{tang}\tfrac{1}{2}U=\mathrm{tang}\tfrac{1}{2}w\mathord{.}(\frac{1+\varepsilon\cos w}{1-\varepsilon\cos w})^{\frac{1}{2}\varepsilon}
\]
\noindent Das ergibt eine für die höhere Geodäsie sehr brauchbare Transformation, deren Verwendung wir hier jedoch nur in wenigen Punkten kurz andeuten können. Wir betrachten die Punkte auf der Oberfläche des Ellipsoids und der Kugel als einander entsprechend, wenn sie dieselbe Länge haben und ihre Breiten \(90^0-w\) beziehungsweise \(90^0-U\) durch die angegebene Gleichung zusammenhängen. Dann entspricht einem System verhältnismäßig kleiner Dreiecke, die auf der Oberfläche des Sphäroids von kürzesten Linien gebildet werden, ein System von Dreiecken auf der Kugeloberfläche. Die Dreiecke, die für die wirkliche Messung infrage kommen, sind stets verhältnismäßig klein. Die Winkel auf der Kugel sind den entsprechenden Winkeln auf dem Sphäroid \textit{genau} gleich. Die Seiten weichen von Großkreisbögen so wenig ab, dass sie in den meisten Fällen, in denen keine äußerste Genauigkeit verlangt wird, als mit ihnen zusammenfallend betrachtet werden können. Wo die höchste Genauigkeit erforderlich ist, lässt sich die Abweichung vom Großkreis mit einfachen Formeln leicht in der nötigen Genauigkeit berechnen. Nachdem man eine Dreiecksseite auf die Kugeloberfläche richtig übertragen hat, kann man daher das ganze System anhand der Winkel so berechnen, als läge es auf der Kugel selbst, nötigenfalls mit der eben angedeuteten Anpassung. Für alle Punkte des Systems kann man die Werte von \(T\) und \(U\) bestimmen und von den letzteren auf die entsprechenden Werte von \(w\) zurückgehen, am einfachsten mit einer sehr leicht zu erstellenden Hilfstafel.
